\section{Related Work Before FLP}

FLP 论文讨论的是异步消息传递系统中的共识不可能性。这个问题并不是孤立出现的；在 FLP 之前，分布式系统、容错计算和分布式数据库领域已经围绕 agreement、commit 和 fault tolerance 形成了多条研究脉络。FLP 的贡献可以放在这些工作之间理解：一方面，它继承了早期 agreement 问题对一致性的关注；另一方面，它把模型收窄到可靠消息、crash failure 和完全异步执行，从而揭示了一个不同于 Byzantine fault 的困难来源。

早期最重要的一条相关研究是 Byzantine agreement。Pease、Shostak 和 Lamport 研究了存在故障处理器时，非故障处理器如何对信息形成一致认识 \cite{pease1980reaching}。他们使用 interactive consistency 的形式化表述，要求所有非故障处理器最终得到相同的信息向量，并且向量中对应非故障处理器的值必须是该处理器真实持有的值。该工作说明，在故障处理器可能向不同接收者发送不同信息时，agreement 问题不能依靠简单多数投票解决，而需要足够的冗余。Lamport、Shostak 和 Pease 随后用 Byzantine Generals problem 将这一类问题表述为经典模型 \cite{lamport1982byzantine}。这些工作关注的是任意或恶意故障，而 FLP 关注的是更弱的 crash failure；这种对比使 FLP 的结论更突出：即使故障进程不会撒谎、不会伪造消息，只是静默停止，完全异步模型仍然会带来根本困难。

另一条与 FLP 动机直接相关的研究线来自分布式数据库中的事务提交。事务提交要求多个数据管理进程一致地决定 commit 或 abort；如果不同副本作出不同决定，数据库状态就会分裂。Dolev 和 Strong 研究了 distributed commit with bounded waiting，试图避免 commit 协议中某些参与者无限等待的问题 \cite{dolev1982distributed}。Skeen 的 nonblocking commit protocols 和 decentralized termination protocol 进一步分析了提交协议中哪些状态会导致 blocking，以及在站点故障后如何推动事务进入最终决定 \cite{skeen1981nonblocking,skeen1982decentralized}。Skeen 和 Stonebraker 则从 crash recovery 角度形式化分布式系统在故障和恢复后的状态一致性问题 \cite{skeen1983formal}。这些工作表明，FLP 论文中的 transaction commit 例子不是附带说明，而是当时分布式数据库系统面临的实际问题：系统既要保持一致性，又希望在故障出现后不无限等待。
